\subsection{高分母 support の現在の必要条件}
\label{sf14:section}

\begin{corollary}[残る support と unicyclic の一型]
\label{sf14:remaining-support}
正の局所整数 charts，marked clutching，対向 $D$ の相殺，node 非負性を満たす
五次の標準 $K_5$ branch graph を考える．高分母 support $H$ が非空なら，
辺を持つ成分は一つだけで，四または五頂点を使い，少なくとも五辺を持つ．
また，少なくとも一頂点の次数は三以上である．
この成分が unicyclic なら，三角形の異なる二頂点に葉辺を一本ずつ付けた graph に限られる．
これらの必要条件と本節までの graph 障害を通過する型は，標識を忘れると十五，
$K_5$ の標識を保つと四百六十一である．これらの存在を主張するものではない．
\end{corollary}

\begin{proof}
成分数と頂点数の条件は系~\ref{fo13:remaining-support}による．
四辺以下なら，四頂点の三角形に葉辺を一本付けた graph だけが新たな候補となるが，
その三角形には二つの二価頂点があるので，定理~\ref{tn14:no-degree-two-triangle}に反する．
従って辺数は五以上である．

五頂点以下の connected unicyclic graph は，唯一の閉路に木を付けた形である．
五角形と四角形成分は排除済みであり，四角形に葉辺を付けた形も
定理~\ref{cp14:no-pendant-square}に反する．
三角形に長さ二の道を付ける場合，または同じ頂点に二本の葉辺を付ける場合には，
三角形の二頂点が二価のままなので，定理~\ref{tn14:no-degree-two-triangle}に反する．
残るのは，三角形の異なる二頂点に葉辺を一本ずつ付けた場合だけである．

最後の型数は，$K_5$ の全 $2^{10}=1024$ 部分 graph から，木成分・単純閉路成分，
二価二頂点を持つ三角形，および葉辺付き四角形を除いて得られる．
全列挙と置換による同型判定は
\nolinkurl{verify_c4_pendant_v14.py} に保存した．
各型の代表と今後の入口は研究ノートに記載する．
\end{proof}
